\scanpage{157}
% archive foliation: 156
\uline{Proposition 3.1.} (\unsure{Remonter}) $X$ étale sur $Y$ (\unsure{par ex. fini}).
$G$ groupe fini \deleted{de}\inserted{opérant par} $Y$-automorphismes de $X$.

\inserted{(i)} \deleted{Si $s\in$} L'ens des pts de $X$ où $G$ opère \inserted{\unsure{sans}} ramification est à la fois ouvert et fermé.
\margin{\unsure{sans inertie}}

(ii) \deleted{$\unsure{X/G\to X}$} $X/G$ existe, et $X\to X/G$ et $X/G-Y$ sont \uline{étales}.

\figure{\begin{tikzpicture}[x=1cm,y=1cm]
\draw (0,2) -- (0,1);
\draw[->] (.15,.7) .. controls (-.6,.1) and (-.25,.3) .. (.8,1);
\draw (0,0) -- (0,-.6);
\end{tikzpicture}}

Si $s\in G$ \deleted{\unsure{opère trivialt}} est dans le groupe d'inertie d'un pt $x\in X$, il \deleted{laisse fixe}\inserted{opère triv. sur tte} la composante connexe de $x$ (vérif), donc l'ens $X(s)$ des $x\in X$ tels que $s\in G_i(x)$ est une réunion de composantes connexes, donc ouvert et fermé.
Comme \deleted{$\unsure{X}$} l'ens des pts à $G_i(x)\ne e$ est la réunion des $X(s)$ $(s\ne e)$, (i) s'ensuit.
D'autre part, $X$ étant qu. affine sur $Y$, $X/G$ existe [---]
\deleted{donc} \unsure{on a} \unsure{finissons-en pour $(ii)$}.

Soit $I$ l'ens des comp. connexes de $X$.
$G$ opère sur $I$, \deleted{\unsure{et on est ramené à}}\inserted{et $X$ est réunion des}
\unsure{un|ens.} $\deleted{\unsure{X(\alpha)}}$ $\alpha\in I/G$, détermine un ouvert $X(\alpha)=\bigcup_{i\in\alpha}X_i$, stable sous $G$.
Donc $X/G=\bigcup X(\alpha)/G$ on est
